\documentclass[conference,a4paper]{IEEEtran}
\usepackage[T1]{fontenc}
\usepackage[utf8]{inputenc}
\usepackage{graphicx,booktabs,amsmath,url}
\begin{document}
\title{Accuracy Budgets and Routing Overhead in Green Edge AI for Environmental Sensor Networks}
\author{\IEEEauthorblockN{Abeer Saleh Alshammari}\IEEEauthorblockA{Independent Researcher\\Kuwait\\engineera209@gmail.com}}
\maketitle
\begin{abstract}Selective inference is attractive for sensor networks because a small model can handle routine inputs while a larger model processes difficult cases. However, avoiding larger-model calls does not necessarily reduce end-to-end computation. This paper evaluates a transparent volatility gate for one-hour-ahead PM2.5 prediction using 409,180 eligible observations from 12 Beijing monitoring stations. Models are trained on 2013--2014 data, routing thresholds are selected on 2015 data, and 118,868 observations from 2016--2017 form the held-out test set. Under a 2\% validation mean absolute error allowance, routing avoids 17.22\% of boosting calls, but median batch inference is 689.98 ms compared with 664.03 ms for always-on boosting. Under a 5\% allowance, validation selects ridge regression alone; its test error is 4.97\% higher and inference is 14.86 times faster than boosting. Ridge also has lower test root mean squared error. A matched-rate random-routing control and station-level analysis expose the limitations of the volatility heuristic. The study provides a reproducible cost--accuracy assessment and an analytical energy break-even condition, without equating invocation reduction with measured energy savings.\end{abstract}
\begin{IEEEkeywords}Green AI, edge inference, sensor networks, adaptive computation, accuracy budget, environmental monitoring\end{IEEEkeywords}
\section{Introduction}
Environmental monitoring generates repeated prediction workloads at gateways serving many sensor streams. A gateway can forecast the next reading locally, support early warning, or forward a prediction to an upstream application. B5G-enabled monitoring motivates efficient analytics, but selecting a machine-learning policy remains an independent system decision. A more accurate predictor may require additional computation, while a small predictor may be sufficient under an explicitly stated error budget. The relevant question is therefore whether conditional computation improves the complete cost--accuracy trade-off after routing overhead is included.

Green AI advocates evaluating efficiency alongside predictive quality [1]. An adaptive policy appears to offer an intuitive compromise: invoke a larger predictor only when the input appears difficult. Yet a reduction in invocation count is only a workload proxy. The small predictor, routing calculation, allocation of subset arrays, and dispatch to the larger predictor all incur costs. If the retained heavy workload remains large, these costs can offset the apparent saving. A simpler static model can also dominate a cascade when the application tolerates a modest increase in average error.

This study asks three questions. First, how does an error allowance chosen before test evaluation affect the fraction of observations sent to a larger model? Second, do avoided calls correspond to shorter measured inference time? Third, does a recent-change heuristic allocate the larger model more effectively than random routing at the same average rate? The contribution is an empirical audit of a deliberately simple routing baseline, rather than a claim of a new general-purpose inference architecture.

We evaluate public environmental observations under chronological separation, report both mean absolute error (MAE) and root mean squared error (RMSE), and distinguish measured runtime from hypothetical energy accounting. The resulting negative and positive findings are both operationally useful: a strict budget can leave too little removable work to offset overhead, whereas a relaxed budget can make an adaptive mechanism unnecessary.

\section{Related Work and Scope}
BranchyNet introduces intermediate exits that allow sufficiently confident predictions to terminate before all neural-network layers are evaluated [2]. IDK Cascades composes pretrained predictors to reduce inference cost [3]. These studies establish conditional computation as an existing approach; our two-regressor switch is a basic instance of that family, not a new cascade concept.

Predictive Exit uses a lightweight prediction mechanism to anticipate exit locations and adjust computation [4]. QUTE addresses resource-efficient uncertainty estimation and model monitoring for TinyML systems [5]. Both illustrate why routing or monitoring costs should be considered. Unlike these neural architectures, the present benchmark uses a linear regressor and histogram gradient boosting, with a gate calculated directly from recent sensor changes. Its value is transparency and easy reproduction.

The study does not benchmark a B5G radio, simulate a communication protocol, or measure network energy. It evaluates an analytics component that could operate at a sensor-network gateway. No claim is made that predicting pollution is itself Green AI: the Green AI question concerns the computational policy used for prediction. Similarly, no state-of-the-art forecasting claim is made because deep temporal models and broad hyperparameter optimization are outside the experiment.

\section{Dataset and Experimental Protocol}
The Beijing Multi-Site Air Quality dataset contains hourly measurements from 12 monitoring stations over March 2013 through February 2017 [6]. It is distributed through the UCI Machine Learning Repository under CC BY 4.0. We downloaded the archive directly from UCI and retained its checksum in the reproducibility metadata. The original archive contains 420,768 station-hour records. After requiring observed current and next-hour PM2.5 values and an exactly one-hour target interval, 409,180 pairs remain. Missing predictor values are retained rather than filled from future observations.

The target is PM2.5 concentration at time t+1. Inputs available at t include current PM2.5; PM2.5 at t-1, t-2, and t-24; current PM10, SO2, NO2, CO, O3, temperature, pressure, dew point, rain, and wind speed; sine and cosine encodings of hour and month; and station and wind-direction indicators. The resulting design has 46 columns. Wind-direction missingness yields an all-zero direction encoding. Lag features are constructed within station after timestamp sorting. No backward interpolation or future weather value is used.

Partitions are assigned using target timestamps. Training targets precede 1 January 2015, validation targets fall in calendar year 2015, and test targets start on 1 January 2016. The respective sizes are 188,119, 102,193, and 118,868. A target at a partition boundary can use earlier historical input, as in chronological deployment. Test inputs can use previously observed test-period readings, but test targets are never used for fitting or threshold selection. The design tests later observations from known stations; it does not test transfer to unseen sites.

The lightweight model is ridge regression with penalty 10, preceded by training-fitted median imputation and standardization. The larger model is scikit-learn histogram gradient boosting with squared-error loss, 150 boosting iterations, 31 maximum leaves, learning rate 0.08, L2 penalty 1, random seed 42, and early stopping disabled [7]. Boosting handles missing numerical values internally. Hyperparameters are fixed for this compact benchmark; neither model is selected through test performance. Persistence, which predicts the current PM2.5 value, provides a no-training baseline.

\begin{table}[t]\caption{Chronological target partitions}\centering\begin{tabular}{lll}\toprule Partition & Target period & Pairs\\\midrule Training & 2013--2014 & 188,119\\Validation & 2015 & 102,193\\Test & 2016--2017 & 118,868\\\bottomrule\end{tabular}\end{table}
\section{Routing and Cost Evaluation}
For station s, the routing score is z(s,t)=|PM2.5(s,t)-PM2.5(s,t-1)|. If the previous reading is missing, z is set to positive infinity. The gate invokes boosting when z exceeds threshold $\tau$; otherwise it returns the ridge prediction. A volatile input is therefore treated as potentially requiring additional computation, although volatility is not a calibrated estimate of relative model advantage.

Candidate thresholds consist of the 101 empirical quantiles of finite validation scores, zero, and the two infinite endpoints. Duplicate thresholds are removed. For each relative allowance $\epsilon$ in {0, 0.02, 0.05, 0.10}, we choose the candidate with the smallest heavy-model invocation rate whose validation MAE is at most (1+$\epsilon$) times the validation MAE of always-on boosting. Ties are broken by validation MAE. Including the all-heavy endpoint ensures feasibility. This is a validation constraint, not a mathematical guarantee on future error.

Each cascade prediction is evaluated using precomputed predictions for accuracy, then separately verified against actual conditional batch execution. The timed cascade first computes ridge predictions, constructs a mask, and calls boosting only for selected rows. This implementation intentionally includes redundant ridge work for selected observations, subset copying, and library dispatch. It represents a simple deployable implementation; a fused or gate-first implementation could improve performance and is not evaluated here.

Inference times cover one full held-out batch with preconstructed features in memory. One warm-up is followed by nine repetitions; the median and interquartile range are reported. CPU numerical libraries are limited to one thread. The platform is Linux x86\_64 with Python 3.12.14 and scikit-learn 1.8.0. Feature construction, data loading, radio transmission, and training are excluded from inference timing. Measurements are batch throughput observations, not per-message latency on a microcontroller. Sequential method timing on a shared host also limits the precision of small differences.

For an allocation control, 30 independently seeded random gates use the observed test invocation probability of each selected policy. Their mean MAE and across-seed standard deviation quantify random allocation variability, not independent repetitions of model training. Error uncertainty is assessed with 2,000 paired bootstrap resamples of calendar-day clusters, keeping all stations and observations from a sampled day together. Daily clustering partially respects dependence; longer cross-day correlation can still make intervals optimistic.

Energy is considered only analytically. Let EH be energy per heavy prediction, EL lightweight prediction energy, and EG routing overhead. If q is the heavy invocation fraction, a simple additive model gives Eroute/EH=$\rho$+q, where $\rho$=(EL+EG)/EH. Break-even requires $\rho$<1-q. This model assumes constant per-item costs and excludes idle energy and network transfers. Neither energy in joules nor carbon emissions is measured or inferred directly from runtime.

\begin{equation}\hat{y}_t=\begin{cases}f_H(x_t),& z_t>\tau,\\f_L(x_t),&z_t\leq\tau.\end{cases}\end{equation}
\begin{equation}\frac{E_{\mathrm{route}}}{E_H}=\rho+q,\qquad \rho<1-q\;\text{for a reduction}.\end{equation}
\section{Results}
Always-on boosting obtains test MAE 9.7485 $\mu$g/m$^3$, compared with 10.0717 for persistence and 10.2333 for ridge. However, ridge has lower RMSE: 19.1527 versus 20.1307 for boosting and 20.0481 for persistence. Model ranking therefore depends on whether the application prioritizes average absolute error or larger errors. Boosting improves MAE over persistence by approximately 3.21\%, a modest advantage that must be weighed against additional computation.

At $\epsilon$=0, validation selects the all-heavy endpoint and the cascade reproduces boosting predictions. Its 803.86 ms median inference time exceeds the 664.03 ms direct boosting call because the wrapper also executes ridge and routing. At $\epsilon$=0.02, the selected threshold is 1 $\mu$g/m$^3$. The validation heavy fraction is 83.24\% and validation MAE is 10.2556, compared with 10.1224 for always-on boosting. On test data the heavy fraction is 82.78\%, so 17.22\% of calls are avoided. Test MAE rises by 0.1537 $\mu$g/m$^3$, or 1.58\%, to 9.9022.

The 2\% cascade takes 689.98 ms per test batch, with interquartile range 666.43--713.18 ms. Always-on boosting takes 664.03 ms, with interquartile range 627.53--677.94 ms. Thus the cascade shows no observed median speedup despite avoiding calls; its median is approximately 3.91\% slower. The overlapping timing ranges and uncontrolled shared-host activity preclude a strong statistical claim about a small runtime difference. The supported conclusion is that invocation reduction did not demonstrate an end-to-end throughput benefit in this implementation.

At $\epsilon$=0.05 and $\epsilon$=0.10, validation selects the all-ridge endpoint. Ridge validation MAE is 10.5406, only 4.13\% above the boosting validation baseline. Test MAE is 4.97\% above boosting, while standalone ridge takes 44.69 ms, yielding a 14.86-fold batch speedup. A cascade wrapper at the all-ridge endpoint takes 49.45 ms in the 5\% configuration; static ridge deployment avoids that extra wrapper cost. The serialized ridge pipeline occupies 2,936 bytes versus 563,485 bytes for boosting. These are pickle sizes, not deployed memory or peak resident memory measurements.

The matched-rate random gate has mean test MAE 9.8315 with standard deviation 0.0099 across 30 seeds at the 2\% policy invocation rate. Its lower mean error than volatility routing indicates that recent change is not an effective relative-benefit signal in this experiment. This comparison is descriptive: the random control is matched to the realized test call rate and is not a separately validation-selected deployment policy. It nevertheless tests whether the heuristic allocation appears better than an uninformative allocation at comparable expected work.

The paired daily-cluster bootstrap interval for the 2\% cascade MAE increase over boosting is [0.1280, 0.1842] $\mu$g/m$^3$. The increase is positive at all 12 stations, ranging from approximately 0.093 at Tiantan to 0.237 at Gucheng. For ridge, the pooled MAE increase interval is [0.2641, 0.6842] $\mu$g/m$^3$. These intervals describe this fixed chronological test period and fitted models; they do not quantify across-city or across-hardware generalization.

A descriptive high-target subgroup further explains the RMSE difference. For 16,282 observations with target PM2.5 at least 150 $\mu$g/m$^3$, ridge MAE is 20.879, persistence MAE is 21.232, and boosting MAE is 23.055. On the remaining 102,586 observations, boosting MAE is 7.637 versus ridge MAE 8.544. This subgroup was inspected after the primary experiment, uses future target values only for analysis, and is not a routing feature or a prospective policy guarantee.

With q=0.8278, the analytical energy model requires $\rho$<0.1722 for positive energy reduction. Assumed $\rho$ values of 0.05, 0.10, and 0.20 give modeled reductions of 12.22\%, 7.22\%, and -2.78\%, respectively. These are sensitivity scenarios, not measured energy savings. They demonstrate why reporting 17.22\% avoided calls as 17.22\% energy reduction would be unjustified.

\begin{table}[t]\caption{Held-out accuracy and median batch runtime}\centering\footnotesize\begin{tabular}{lrrrr}\toprule Policy & MAE & RMSE & Heavy \% & ms\\\midrule Persistence & 10.072 & 20.048 & 0 & 0.57\\Ridge & 10.233 & 19.153 & 0 & 44.69\\Boosting & 9.748 & 20.131 & 100 & 664.03\\Gate 0\% & 9.748 & 20.131 & 100 & 803.86\\Gate 2\% & 9.902 & 20.166 & 82.78 & 689.98\\Gate 5\% & 10.233 & 19.153 & 0 & 49.45\\Gate 10\% & 10.233 & 19.153 & 0 & 45.00\\\bottomrule\end{tabular}\\[2pt]\parbox{.97\columnwidth}{MAE and RMSE in $\mu$g/m$^3$. Gate percentages denote validation error allowances, not guaranteed test errors. Persistence timing includes array extraction.}\end{table}
\begin{figure}[t]\centering\includegraphics[width=\columnwidth]{runtime.pdf}\caption{Measured runtime and error for the fixed test batch. Conditional routing does not establish a speedup over boosting.}\end{figure}
\begin{figure}[t]\centering\includegraphics[width=\columnwidth]{stations.pdf}\caption{The 2\% cascade increases MAE at every station relative to boosting.}\end{figure}
\begin{figure}[t]\centering\includegraphics[width=\columnwidth]{energy.pdf}\caption{Analytical energy sensitivity for $q=0.8278$. The dashed line marks break-even. Values are assumptions, not energy measurements.}\end{figure}
\section{Implications and Limitations}
The first practical implication is that an allowable error budget should be specified before choosing adaptive inference. Under a tight budget, a simple gate can retain most of the expensive workload and add overhead. Under a looser budget, a static small model can be the more efficient choice. The validation-selected all-ridge endpoint is therefore a meaningful deployment result, not evidence that the adaptive gate itself succeeded.

The second implication concerns the loss function. Ridge loses to boosting on aggregate MAE but wins on RMSE and the descriptive high-target subgroup. A monitoring service emphasizing large pollution excursions should not select a policy from mean MAE alone. Future evaluation should specify an operational alert cost or asymmetric loss before calibration and assess whether any runtime saving remains under that objective.

Important limitations remain. Only one historical dataset and one train--validation--test partition are used. The stations provide environmental observations rather than a live IoT gateway trace. No B5G link, sensor power, communication payload, hardware energy, or carbon intensity is measured. Both models require the same feature availability; lightweight inference does not imply fewer sensing operations or less uplink traffic. Maintaining two models can increase memory and training costs. Nine timings do not establish hardware-independent speedups.

The volatility gate is a transparent baseline with weak allocation results. Future work should compare a benefit-prediction gate, a gate-first implementation, temporally blocked calibration at multiple origins, and longer-block uncertainty estimates. Hardware-in-the-loop measurements should jointly capture per-record latency, training amortization, idle power, and actual energy. Such extensions are necessary before claiming an energy-efficient operational B5G system. The present study supports a narrower conclusion about accuracy budgets, implementation overhead, and static alternatives.

\section{Conclusion}
A chronological benchmark of 12 environmental monitoring stations shows that avoided larger-model calls are insufficient evidence of Green AI benefit. A validation-selected 2\% error-allowance cascade avoids 17.22\% of boosting invocations but provides no observed batch speedup and allocates computation less effectively than a matched-rate random control. With a 5\% validation allowance, the selected policy collapses to ridge alone, yielding a 14.86-fold measured batch speedup with a 4.97\% test MAE increase and lower RMSE. Green edge analytics should report the complete cost--accuracy trade-off, include static baselines, and separate measured computation from modeled energy.

\section*{Acknowledgment}
OpenAI ChatGPT with Codex was used to develop the experimental Python code, generate plots from computed outputs, and draft Sections I--VII and the abstract. The experiments were executed on the accompanying dataset, and numerical statements were derived from their saved outputs. No synthetic observations or invented experimental results were used. Author review of the manuscript, code, references, and scientific claims is required before submission. Claude was not used in producing this version.

\begin{thebibliography}{7}
\bibitem{r1} R. Schwartz, J. Dodge, N. A. Smith, and O. Etzioni, ``Green AI,'' Communications of the ACM, vol. 63, no. 12, pp. 54--63, 2020, doi: 10.1145/3381831.
\bibitem{r2} S. Teerapittayanon, B. McDanel, and H. T. Kung, ``BranchyNet: Fast inference via early exiting from deep neural networks,'' arXiv:1709.01686, 2017.
\bibitem{r3} X. Wang, Y. Luo, D. Crankshaw, A. Tumanov, F. Yu, and J. E. Gonzalez, ``IDK cascades: Fast deep learning by learning not to overthink,'' arXiv:1706.00885, revised 2018.
\bibitem{r4} X. Li et al., ``Predictive Exit: Prediction of fine-grained early exits for computation- and energy-efficient inference,'' Proceedings of the AAAI Conference on Artificial Intelligence, vol. 37, no. 7, pp. 8657--8665, 2023, doi: 10.1609/aaai.v37i7.26042.
\bibitem{r5} N. P. Ghanathe and S. J. E. Wilton, ``QUTE: Quantifying uncertainty in TinyML models with early-exit-assisted ensembles for model-monitoring,'' Proceedings of Machine Learning Research, vol. 267, pp. 19286--19306, 2025.
\bibitem{r6} S. Chen, ``Beijing Multi-Site Air Quality,'' UCI Machine Learning Repository, 2017, doi: 10.24432/C5RK5G.
\bibitem{r7} Scikit-learn developers, ``HistGradientBoostingRegressor,'' scikit-learn 1.8.0 documentation. [Online]. Available: \url{https://scikit-learn.org/1.8/modules/generated/sklearn.ensemble.HistGradientBoostingRegressor.html}. Accessed: Oct. 30, 2026.
\end{thebibliography}
\end{document}
